\section{Fazit und Ausblick}

Die vorliegende Arbeit zeigt, dass die Entwicklung der Perspektive weit mehr war als die rein geometrische Anordnung von Objekten im Raum. Das Streben der Künstler und Mathematiker, Anatomie, Licht und Emotionen so lebensecht einzufangen - man denke hier nur an Vorzeigekünstler wie Leonardo Da Vinci -  als könne die Person jeden Moment aus dem Rahmen steigen, bildet heute die Grundlage für völlig neue Forschungszweige. So ermöglicht diese historische Detailtreue modernen Experten wie dem Pathologen Prof. Dr. Andreas Nerlich der Ludwig-Maximilians-Universität München die gezielte Analyse krankheitstypischer Merkmale auf alten Gemälden -- ein Verfahren, das bei flachen Darstellungen undenkbar wäre. Ein Beispiel ist Leonardo da Vincis \textit{Mona Lisa}, in welcher gelbliche Veränderungen am Augenlid (Xanthelasmen) und Lipome festzustellen sind, die auf eine Fettstoffwechselstörung hindeuten. \footnote{Vgl.: \cite{Morgenpost2026}} 

In diesem Sinne öffnete Leon Battista Alberti mit seiner historischen Analogie nicht nur ein Fenster, um die Realität als Momentaufnahme festzuhalten, sondern er spannte einen völlig neuen Raum der Möglichkeiten auf. Dieser Raum erstreckt sich bis in die Gegenwart: Er erlaubt es uns heute, die einst abstrakte Mathematik im Glanze moderner Werkzeuge wie Manim durch lebendige Animationen geradezu "`magisch"' neu zu vermitteln.